\documentclass[12pt]{article}
\usepackage[utf8]{inputenc}
\usepackage{booktabs}
\usepackage{multicol}
\usepackage{graphicx}
\usepackage{fancyhdr}
\usepackage{lipsum}
\usepackage[style=numeric-comp]{biblatex}
\addbibresource{bibliography.bib}
\setlength{\columnsep}{4cm} %a djust base on your name length

\begin{document}
\pagenumbering{gobble}

%% Copertina
\begin{center}
{\LARGE \textbf{Conservatorio di Musica Santa Cecilia}} \\
\vspace{0.2cm}
{\Large {Dipartimento di Nuove Tecnologie e Linguaggi
Musicali}} \\ 
\vspace{1cm}

\includegraphics[width=12cm]{figures/logo-CSC-pdf-cropped.pdf} \\
\vspace{0.8cm}

{\Large {Tesi di Laurea Biennale in Musica Elettronica}} \\
\vspace{1cm}

{\LARGE \textbf{Sistemi Complessi Adattivi per la performance musicale in Live Electronics}} \\ 
\vspace{1cm}

\end{center}

\begin{multicols}{2}
\noindent \large{Relatore:} \\
\large{\textbf{Giuseppe Silvi}} \\
\vspace{0.1cm}

\noindent {\large {Correlatori:}} \\
\large{\textbf{Agostino Di Scipio}} \\
\large{\textbf{Dario Sanfilippo}} \\
\vspace{0.1cm}

%% \noindent {\large {Controrelatore:}} \\
%% \large{\textbf{Prof. Name}} \\
%% \vspace{0.1cm}
\columnbreak

\noindent{\large {Candidato:}} \\
\large{\textbf{Luca Spanedda}} \\
\end{multicols}
\vspace{2cm}

\begin{center}
    \large{\textbf{Anno Accademico 2021/2022}}
\end{center}
\newpage
\ 
\newpage

%% Dichiarazione di intenti
\vfill
\LARGE \textbf{Dichiarazione} \normalsize \newline \newline
Dichiaro che il sottoscritto nonché autore del
documento è il responsabile del suo contenuto, 
e per le parti tratte da altri lavori, 
queste vengono espressamente dichiarate
citando le fonti.
\
\newline
Luca Spanedda.
\ 
\newpage

%% Ringraziamenti
\vfill
\LARGE \textbf{Ringraziamenti} \normalsize \newline \newline
qui i ringraziamenti.
\newpage
\
\newpage

%% Abstract
\begin{abstract}
Il lavoro presentato qui è uno studio di implementazione, analisi ed esecuzione
di tre Sistemi Complessi Adattivi (CASs: Complex Adaptive Systems) 
creati appositiamente per la performance musicale (in Live Electronics),
la scelta di questi tre sistemi corrisponde a tre diversi casi di studio 
(case study) nell'implementazione di dinamiche nonlineari sfruttate 
per la generazione dei comportamenti emergenti nei Sisitemi Complessi.
Una prima parte del lavoro tratterà dell'implementazione e l'analisi di due brani,
creati rispettivamente da due compositori e ricercatori riconosciuti 
nell'ambito internazionale della Computer Music: Agostino Di Scipio e Dario Sanfilippo.
Di Agostino di Scipio studieremo un sistema con nonlinearità provenenti dal dominio analogico, 
che sfrutta fenomeni generativi nel mondo fisico riportati all'interno del sistema 
grazie alle conversioni AD/DA (analogico a digitale e viceversa). 
Di Dario Sanfilippo studieremo invece un sistema che sfrutta nonlinearità 
appositamente programmate dal compositore nel mondo digitale DSP (Digital Signal Processing),
tramite alcuni principi di autoregolazione, 
senza interazioni e perturbazioni provenienti dal mondo esterno. 
Infine la seconda parte del lavoro tratta l'implementazione di un mio brano,
ibrido fra i due casi di studio presentati qui,
che andrà a conclusione del lavoro di ricerca, analisi ed implementazione
svolto durante il corso della tesi.
\end{abstract}
\newpage
\

% conteggio pagine
\pagenumbering{arabic}

\tableofcontents
\listoftables
\listoffigures
\newpage

\pagestyle{fancy}


%% --------------------------------------------------- Capitoli ----------------
%% capitolo TEST
%%\section{Introduction}
\label{sec:introduction}

% write your text here..
A citation~\cite{Hasselmo5249}. 
\lipsum[1-3]

\subsection{A subsection}
\label{sec:asubsection}

% write your text here..
\lipsum[1-2]

\begin{figure}
    \centering
    \includegraphics[width=0.25\textwidth]{figures/unito-logo.png}
    \caption{A figure}
    \label{fig:figure}
\end{figure}

\lipsum[3-4]

\begin{table}
  \caption{Sample table}
  \label{sample-table}
  \centering
  \begin{tabular}{lll}
    \toprule
    \multicolumn{2}{c}{Part}                   \\
    \cmidrule(r){1-2}
    Name     & Description     & Size ($\mu$m) \\
    \midrule
    Dendrite & Input terminal  & $\sim$100     \\
    Axon     & Output terminal & $\sim$10      \\
    Soma     & Cell body       & up to $10^6$  \\
    \bottomrule
  \end{tabular}
\end{table}

\subsubsection{A sub-subsection}

\lipsum[5]

\paragraph{A paragraph} \lipsum[6]

%% TITOLO
\section{Introduzione}
\label{sec:Introduzione}

\begin{center}
\vspace{0.5cm}
\textit{Per primo fu Caos}
\\
Esiodo - Teogonia
\vspace{0.5cm}
\end{center}

Nel mito greco il principio del mondo non è opera di una divinità creatrice,
ma ogni cosa deriva dalla materia, dall’evoluzione del caos primordiale,
e più in generale nelle mitologie antiche, il caos è quasi sempre
contrapposto al cosmo, nel senso di universo disordinato il primo e ordinato
il secondo. \\
I concetti di caos e di ordine strutturato per come sono utilizzati dalla scienza moderna, fino alla fine del XIX secolo non appartenevano alla fisica ed erano considerati argomento “intrattabile”.
Il concetto di ordine permeava tutta la meccanica,
ma nella forma semplice di un solo tipo di ordine: l’ordine della ripetizioni periodica.
La Terra che ruota intorno al Sole o il pendolo ideale che oscilla in eterno, ne sono i classici esempi. \\
La situazione cambiò notevolmente nel Novecento:
dopo la crisi dei fondamenti delle scienze di inizio secolo,
le Scienze della Complessità e la Sistemica modificano radicalmente l’ottica da cui guardare questi problemi. [\cite{caosvillani}]
Al contempo le nascenti considerazioni strutturali e teoriche
nella musica alla ricerca di vie al di fuori del sistema tonale,
parallelamente all'esigenza di introdurre nuovi paradigmi all'interno delle scienze,
hanno contribuito all'avvenire di importanti
punti di incontro fra i due ambiti. \\

\subsection{La Cibernetica}
\label{sec:La Cibernetica}
Uno dei più importanti avanzamenti nelle scienze
al termine della seconda guerra mondiale, risedette nell'introduzione della
cibernetica e della teoria generale dei sistemi, che hanno conseguentemente
portato alla nascita del pensiero sistemico e del concetto
di scienze della complessità. \\
La cibernetica è la scienza che studia i principi astratti di organizzazione
nei sistemi complessi, ed ebbe inizio durante gli anni della seconda guerra
mondiale, merito del fisico e matematico Norbert Weiner.
Nel 1940 Wiener insieme ad altre ad altre prominenti figure provenienti
da diversi ambiti scientifici,
come Ross Ashby, Margaret Mead, Gregory Bateson, Heinz von Foerster,
partecipano ad una serie di conferenze
multidisciplinari chiamate "The Macy Conferences", inizialmente intitolate come
"Feedback Mechanism in Biology and the Social Sciences"
con l'obiettivo comune di andare a definire
gli ambiti di interesse della nuova scienza.
Il concetto sviluppato dai greci - kybernetes -
venne poi ripreso da Norbert Weiner nel 1948,
che ispirato dalla meccanica ed i suoi risultati durante la guerra
e contemporaneamente dallo sviluppo della teoria della comunicazione
(o informazione) di Claude Shannon,
e con la volontà di sviluppare una teoria generalizzata dei principi di
organizzazione e controllo nei sistemi emersi durante le conferenze,
pubblicherà un libro nel 1948:
La cibernetica, controllo e comunicazione nell'animale e nella macchina;
in cui definiva l'ambito di interesse e gli obiettivi della nuova disciplina
inaugurando anche l'uso del nuovo termine da lui coniato.
A seguito di questo libro che riscuoterà
un importante successo, le conferenze presero il nome di
"Cybernetics, Circular Causal, and Feedback Mechanism in Biological and Social Systems",
riconoscendo Wiener come la principale figura di spicco della nuova scienza. \\
In particolare come evidenziato fino ad ora dalla sua natura multidisciplinare,
la cibernetica non si interessa di individuare in cosa consistano questi sistemi,
ma più che altro comprenderne il loro funzionamento:
\\ \\
- come usano l'informazione e come la scambiano gli agenti
\\
- come collaborano fra loro in direzione di un obiettivo comune
\\
- come contrastano il rumore nel trasferimento dell'informazione
\\
  e così via...
\\ \\
Le fortunate premesse iniziali della cibernetica risiedevano in una convinzione
da parte di questi scienziati provenienti dai differenti ambiti disciplinari,
che esistesse uno "schema processuale" comune ad organismi viventi e macchine,
rintracciato attraverso una ricerca uniforme garantita da dell'utilizzo di un metodo
"sintetico" e "comportamentale".
Fra gli anni '60 e la metà del '70, grazie agli scienziati
Heinz von Foerster, Margaret Mead, e altri,
si compierà un ulteriore passo fondamentale che porterà
il pensiero sistemico verso il consolidamento in una scienza più concreta
dando vita alla "Cibernetica di secondo ordine",
anche chiamata come "la cibernetica dei sistemi di osservazione",
che consiste nell'applicazione
ricorsiva della cibernetica a se stessa e la pratica riflessiva della cibernetica
secondo tale critica.
La differenza fra cibernetica di primo e secondo ordine risiese nel fatto,
che mentre nel primo periodo lo studioso di cibernetica (di primo ordine)
studiava un sistema da un punto di vista passivo, da quello dell'osservatore
dei comportamenti di un sistema.
Il  cibernetico di secondo ordine lavora ed interviene nel comportamento
e nella costruzione di un sistema complesso,
riconoscendo il sistema come un agente con cui interagire e
riconoscendo esso stesso come agente nell'interazione col sistema. \\
A partire dalle sue importanti premesse,
la cibernetica ha conseguentemente poi avuto un ruolo centrale nello sviluppo di
molti studi scientifici e la nascita
di nuovi ambiti come: l'intelligenza artificiale, la teoria del caos,
la teoria della catastrofe,
la teoria dei controlli, la teoria generale dei sistemi, la robotica,
la psicologia e le scienze sociali,
ecc.
nella mappa di B.Castellani e L.Gerrits riportata per intero nella pagina seguente,
possiamo visualizzare con più precisione e accurtezza
l'insorgere e l'evoluzione di questi paradigmi scientifici, per averne una
visione più completa relativa al loro sviluppo.

\clearpage

\begin{figure}[!h]
    \centering
    \includegraphics[width=1.0\textwidth, angle=0]{figures/complexitymap.pdf}
    \label{fig:figure}
\end{figure}
%% https://www.art-sciencefactory.com/MapLegend.html

\clearpage

\subsection{Le cibernetiche nella musica}
\label{sec:Le cibernetiche nella musica}
All'inizio degli anni '60 in seno alle nascita delle scienze complesse,
l'uso di sistemi di feedback e la rilevanza dei circuiti informativi chiusi
nelle strutture organizzate,
ha goduto di uno slancio popolare anche nel mondo della musica
e più in generale dell'arte.
Uno dei primi nella storia dell'arte ad evocare l'uso della cibernetica
nei propri lavori è stato
Nicolas Schoeffer con il suo ciclo di lavori "spazio-dinamici",
in particolare ha creato la prima installazione ad implementare meccanismi
di auto-regolazione, il CYSP-1 [\cite{SanfilippoValle}],
capace di essere sensibile all'ambiente esterno e a se stesso
grazie ad una serie di tecnologie offerte dalla compagnia Philips (fotocellule e microfoni),
e reagire sonoramente a questi stimoli riproducendo
una serie di registrazioni composte dal compositore francese Pierre Henry,
collaboratore di Pierre Schaeffer ed insieme a lui figura centrale nella nascita della Musique concrète.
Un'altra importante esperienza del periodo iniziale è quella del compositore Roland Kayn,
che dopo essersi avvicinato alla musica elettronica sotto la guida di Herbert Eimert
nello studio di Colonia (1954),
e dopo essersi trasferito a Roma nel 1960, dal 1964 assieme ad Aldo Clementi e Franco Evangelisti
fonda il Gruppo di improvvisazione Nuova Consonanza, [\cite{Kaynbio}]
del quale fece parte sino al 1968,
ed in quel periodo ispirato dalle teorie della cibernetica iniziò a sperimentare
estensivamente con sistemi di autoregolazione basati su feedback loops,
sia come modelli formali per composizioni strumentali che come reti di generatori di segnale analogici.
Tuttavia, a parte casi popolari di deliberate dichiarazioni
formali da parte degli artisti,
come nel caso di Roland Kayn,
non bisogna pensare a questi lavori appena citati (ed altri riportati a seguito),
come ad atti pioneristici che sancisono la nascita della cibernetica in musica,
ma proprio come si dice per la scoperta del fuoco
lo scenario più accurato risiede probabilmente nel fatto che
tanti autori provenienti da diverse parti del mondo, nello stesso periodo
sono stati influenzati e si sono influenzati a vicenda con le stesse idee
provenienti da un interesse condiviso per le teorie cibernetiche di Weiner e delle Macy Conferences.
Si può pensare ad esempio a quelle che sono le esperienze dello studio di Colonia:
nel 1951, Herbert Eimert e Werner Meyer-Eppler persuasero il direttore della NWDR, Hanns Hartmann,
a creare uno Studio per la Musica Elettronica, che Eimert diresse fino al 1962.
Questo è diventato lo studio più influente al mondo durante gli anni 1950 e 1960,
con ospiti alcuni dei più importanti compositori contemporanei provenienti da tutta europa,
come il già citato Roland Kayn, Franco Evangelisti, Karlheinz Stockhausen, Herbert Brun,
Cornelius Cardew, e molti altri. [\cite{EMSColonia}]
Non a caso in quel periodo il lavoro di ricerca condotto da Werner-Meyer Eppler,
scienziato, musicista ideatore e direttore dello studio di Colonia,
ha trovato sin dalla nascita dello studio fondamento in quelle che sono state
le teorizzazioni della Cibernetica.
C'è poi il caso di Franco Evangelisti,
come citato prima fondatore insieme a Roland Kayn del Gruppo Improvvisazione Nuova Consonanza,
che in quel periodo (qualche anno prima della fondazione del Gruppo a Roma)
si trova nello studio di Colonia per lavorare al brano "Incontri di Fasce Sonore",
e quando farà ritorno a Roma poi citerà più volte deliberatamente in interviste, scritti,
e altre documentazioni, il suo approccio sistemico/cibernetico a quelle che saranno le esperienze con il Gruppo.
Se cambiamo paese e passiamo dall'Europa od osservare l'America in quel periodo,
possiamo pensare a quelli che sono i lavori di Louis e Bebe Barron,
con i circuiti in retroazione destinati al corto circuito
e utilizzati appositamente come materiale per la generazione acustica di trame incise su nastro,
o ai lavori pioneristici di
John Cage, David Tudor, Robert Ashley e Steve Reich,
che sfruttano ed esplorano l'effetto Larsen in modo artistico insieme
alle logiche di Feedback. \\
Un secondo periodo costituito da un approccio sistemico più consapevole
che inizia a tracciare la strada per un pensiero ecosistemico della composizione,
inizia invece dal lavoro
di Alvin Lucier, che nel 1969 scriverà quello che sarà un brano emblematico per
la cibernetica in musica "I'm sitting in a room",
è un altro brano importante per quelle che sono
le logiche di interazione sistemiche fra uomo/macchina/ambiente
e che sancisce una volta per tutte
l'interazione sistemica dove il musicista l'ambiente e lo strumento sono parti di un insieme del sistema "più complesso" con un comportamento collettivo derivato dai singoli agenti,
in un un'interazione con l'ambiente circostante.
In I'm sitting in a room, un performer al centro della stanza
recita in un microfono un testo che descrive il fenomeno che avverrà poco a poco,
la voce recitante nel microfono viene registrata e poi riprodotta da altoparlanti
posti nella stanza, il suono della regitrazione riprodotta da questi altoparlanti
viene registrato nuovamente durante la riproduzione, l'operazione
viene ripetuta in un in una casualità circolare di volta in volta dove alla fine rimarranno
solo i contributi provenienti dalla stanza, dalla voce e dalla catena elettroacustica,
dando vita nel loro insieme ad un effetto Larsen, la natura
nonlineare del processo e degli agenti porterà di volta in volta ad un risutato
sempre differente.
Dopo l'esperienza di Lucier, nel 1974 Nicolas Collins compone "pea soup"
mentre è studente alla Wesleyan University.
Pea soup consiste in una rete adattiva di circuiti analogici (3 Countryman Phase Shifters),
che intona il feedback positivo dell'effetto Larsen ad una frequenza risonante diversa
ogni volta che questo inzia ad emergere.
Ad oggi svariati compositori a partire dalle trame delineate dalle scienze complesse e
dai lavori citati operano nell'ambito della musica elettronica con un approccio sistemico,
un importante caso è quello di Agostino Di Scipio, che contribuisce significatamente
nell'ambito della computer music sin dai primi anni '90, divenendo una
delle figure più importanti nell'area della composizione ecosistemica e nel suo
caso in particolare del live electronics, o quello di Dario Sanfilippo
compositore e ricercatore con all'attivo recenti importanti pubblicazioni e lavori
nell'ambito dei sistemi di feedback in musica, e in particolare di non linearità nei sistemi
in DSP.

\subsection{Il Feedback}
\label{sec:Il Feedback}
Il feedback (o retroazione) è un concetto cibernetico che sta ad indicare
la capacità di un sistema di autoregolarsi tenendo conto degli effetti scaturiti
dalla modificazione delle caratteristiche del sistema stesso.
In termini appartenenti alla fisica, è la capacità di un sistema dinamico
di tenere conto dei risultati del sistema per
modificare le caratteristiche del sistema stesso.
Negli esseri viventi, ad esempio, i sistemi a retroazione negativa e positiva
sono ampiamente utilizzati per regolare l'omeostasi dell'organismo.
Esistono idealmente due tipologie di Feedback:
\\ \\
- a Retroazione Positiva
\\
- a Retroazione Negativa
\\ \\
La retroazione positiva tende ad accelerare un processo,
mentre la retroazione negativa a rallentarlo.
La retroazione negativa aiuta a mantenere la stabilità di un sistema,
contrastando i cambiamenti provenienti dall'ambiente esterno.
Mentre la positiva tende generalmente alla complessità.
Nel controllo di un sistema complesso,
come può essere ad esempio quello del feedback acustico,
introdurre delle linearità tramite retroazione
vuole dire costringere la complessità
a dei comportamenti prevedibili, si può pensare ad esempio
all'intonazione del feedback,
che da un comportamento complesso della sorgente e del ricettore arriva
ad uno lineare.
Mentre introdurre delle non-linearità nel sistema tramite la retroazione,
vuoldire portare questo verso comportamenti non più prevedibili.
Questi due tipi di comportamento possono essere ottenuti per l'appunto
sia velocizzando che rallentando questi processi,
in maniera dipendente dal caso specifico.
I filtri digitali o analogici nell'audio,
possono essere pensati per esempio come uno strumento
di contrasto rispetto a questo tipo di comportamenti:
dove se si allineano le fasi si creano dei poli,
mentre se si disallineano si punta alla complessità del sistema.
Di fatto la storia delle tecnologie elettroacustiche ha in generale
da sempre incorporato il principio
del feedback sin dalle sue origini, basti pensare a tecnologie
come la valvola audion di Lee De Forest (1910),
o i circuiti di feedback negativo di Harold Black (1920) [\cite{echodiscipio}].
Fino al caso del feedback elettroacustico (effetto Larsen), che abbiamo visto
esser stata una delle risorse centrali dei primi compositori cibernetici.
Addentriamoci ora verso una spiegazione più tecnica del feedback acustico
che ci servirà per comprendere più
a fondo la poetica di questi compositori,
parafrasando l'articolo appena esposto di Agostino Di Scipio:
In una stanza vengono collegati fra loro (tramite uno o più stadi di amplificazione) gli
elementi di una catena elettroacustica molto elementare: Microfono ed Altoparlante. Anche in
una situazione di "silenzio" ideale, viene catturata dal microfono in ogni caso,
inevitabilmente, perfino la turbolenza minima appena udibile presente nel rumore ambientale.
Questo suono catturato dal microfono, viene amplificato e riprodotto dall’altoparlante a sua
volta. E se l'amplificazione è sufficiente il suono dall’altoparlante ritorna al microfono e
il design della catena elettroacustica si chiude su se stessa,
creando un circuito di
retroazione (loop), anche detto di feedback.
Il livello di ampiezza, le caratteristiche
tecniche trasduttive di microfono e altoparlante, la loro distanza relativa, la distanza dalle
pareti ed altri potenziali fattori influenti, delineano un’oscillazionein fase con il segnale
che viene a sommarsi ad esso e viene amplificata e riprodotta a sua volta con ampiezza via via
crescente, idealmente illimitata. Con livelli di guadagno non troppo elevati, ciò che si
genera è un fastidio udibile, una sorta di 'alone':
la reiniezione del suono decade più o meno
rapidamente, in una specie di effetto riverbero composto da suono spettralmente irregolare.
Con livelli di guadagno più elevati, il feedback loop entra in un regime di auto-oscillazione.
A causa della reiniezione ripetuta, il rumore di fondo appena udibile ma spettralmente ampio
si accumula nel loop e alla fine (rapidamente) produce un suono sostenuto sempre più forte di
uno spettro più ristretto - questo è spesso sentito come un tono di picco di altezza definita
o un gruppo di toni. E questo è l'effetto Larsen: (dal nome del fisico Søren Absalon Larsen
che per primo ne scoprì il principio), detto anche feedback acustico o più prosaicamente
ritorno, è timbricamente riconoscibile come un tipico fischio stridente, che si sviluppa come
abbiamo detto quando i suoni emessi da un altoparlante vengono captati con sufficiente
"potenza di innesco" da un trasduttore (che può però anche essere oltre al microfono, un
pick-up di uno strumento musicale elettrico, come una chitarra o un basso, o un trasduttore di
altra natura).


\printbibliography
\appendix 
\section{First appendix}

\section{Second appendix}

\end{document}